\chapter{Альтернативные современные теории обучения на основе дофамина}
\chaptermark{Другие теории \nt{DOPA}}
\label{sec:dofaalt}
\begin{chaptergoals}
\item как \nt{DOPA}-нейроны с асимметричным ответом на ошибки записывают всё распределение награды;
\item как часть \nt{DOPA}-нейронов сообщает о важности или опасности, а не о знаке ошибки;
\item как один провод несёт быструю ошибку для обучения и медленный уровень, задающий темп действий;
\item почему сигнал \nt{DOPA} --- жгут, а не одно число, и какая теория оспаривает саму ошибку.
\end{chaptergoals}

\section{Что на самом деле идёт по проводу}

В главе~\ref{sec:dofa} \nt{DOPA}-система была у нас одним проводом, по которому идёт одно число --- ошибка предсказания награды $\delta$~\eqref{eq:ctd}. Эта картина объясняет многое, но чем точнее измеряют дофамин, тем больше находят того, что в неё не укладывается. Одни нейроны отвечают на неприятное, как на награду; концентрация дофамина растёт, пока животное бежит к цели, хотя ничего неожиданного не происходит; разные \nt{DOPA}-нейроны ведут себя по-разному.

Для инженера все эти находки --- вопрос об одном: \emph{какой формат у сигнала на широковещательной шине?} Одно число или вектор? Один сигнал на проводе или несколько, разнесённых по частоте? Говорит ли он о будущем, как $\delta$, или о прошлом? Ниже --- шесть современных ответов; для каждого мы отделим измеренное от предполагаемого.

\section{Не одно число, а распределение}

Ошибка~\eqref{eq:ctd} учит критика \emph{среднему} ожиданию награды. Но верные полкапли сока и капля с вероятностью одна вторая имеют одинаковое среднее. Чтобы их различать, нужно всё распределение награды.

Возьмём простейшую задачу: после предвестника приходит награда случайного размера $r$. Пусть \nt{DOPA}-нейронов много, и каждый, $i$-й, обслуживает свою оценку $V_i$, так что в момент награды его ошибка равна $\delta_i=r-V_i$. И пусть положительную ошибку каждый учитывает с весом $\alpha_i^+$, отрицательную --- с весом $\alpha_i^-$. После каждой награды оценка сдвигается на
\begin{equation}
\Delta V_i \;=\; \eta\,\bigl(\alpha_i^+\,[\delta_i]_+ \;-\; \alpha_i^-\,[-\delta_i]_+\bigr).
\label{eq:asymmetric}
\end{equation}
Оценка перестаёт меняться, когда в среднем положительные поправки уравновешивают отрицательные:
\begin{empheq}[box=\keybox]{equation}
\alpha_i^+\,\bigl\langle[r-V_i]_+\bigr\rangle \;=\; \alpha_i^-\,\bigl\langle[V_i-r]_+\bigr\rangle ,
\qquad
\kappa_i \;=\; \frac{\alpha_i^+}{\alpha_i^++\alpha_i^-} .
\label{eq:expectile}
\end{empheq}
При $\kappa_i=1/2$ это обычное среднее. При $\kappa_i>1/2$ нейрон --- «оптимист»: его оценка сдвинута к верхнему краю распределения, при $\kappa_i<1/2$ --- «пессимист».

Пример: сок приходит с вероятностью $1/2$, размер капли $1$. Тогда~\eqref{eq:expectile} даёт $\kappa_i\cdot\tfrac12(1-V_i)=(1-\kappa_i)\cdot\tfrac12V_i$, то есть $V_i=\kappa_i$ (рис.~\ref{fig:expectiles}). Набор нейронов с разными $\kappa_i$ записывает распределение награды так же, как набор квантилей записывает распределение в статистике.

\begin{figure}[t]
\centering
\begin{tikzpicture}[>=Stealth,x=7cm,y=3cm]
\draw[->,gray!70!black] (-0.08,0) -- (1.15,0) node[right,black] {\small награда};
\draw[->,gray!70!black] (-0.08,0) -- (-0.08,0.75) node[left,black] {\small вероятность};
\fill[blue!25] (-0.03,0) rectangle (0.03,0.5);
\fill[blue!25] (0.97,0) rectangle (1.03,0.5);
\node[below,font=\small] at (0,0) {$0$};
\node[below,font=\small] at (1,0) {$1$};
\node[left,font=\scriptsize] at (-0.08,0.5) {$1/2$};
\foreach \k/\c/\lab/\h in {0.2/blue!60!black/{пессимист, $\kappa=0.2$}/0.62, 0.5/gray!50!black/{среднее, $\kappa=0.5$}/0.42, 0.8/red!70!black/{оптимист, $\kappa=0.8$}/0.22}{
  \draw[very thick,\c] (\k,0) -- (\k,\h);
  \node[above,font=\scriptsize,\c] at (\k,\h) {\lab};}
\end{tikzpicture}
\caption{Распределение награды записано набором \nt{DOPA}-нейронов. Награда $0$ или $1$ с вероятностью $1/2$ (синие столбики); нейрон с долей $\kappa$ приходит к оценке $V=\kappa$~\eqref{eq:expectile}.}
\label{fig:expectiles}
\end{figure}

Что измерено: у разных \nt{DOPA}-нейронов разная «точка переворота» --- размер награды, при котором ответ сменяется со всплеска на провал, --- и нейроны с большей асимметрией ответов на положительные и отрицательные ошибки переворачиваются при больших наградах, как и требует~\eqref{eq:expectile}~\cite{Dabney2020}. Из ответов группы нейронов удаётся восстановить форму распределения. Что это главная функция разброса, а не побочный эффект, --- гипотеза.

\begin{keyblock}
\begin{proposition}[Распределение из асимметрии]
\label{prop:expectile}
Если \nt{DOPA}-нейроны различаются долей $\kappa_i$, с которой они учитывают положительные ошибки, их оценки сходятся к разным точкам~\eqref{eq:expectile} распределения награды. Группа нейронов передаёт не среднее, а распределение, и тем самым --- риск.
\end{proposition}
\end{keyblock}

\section{Не только ошибка, но и важность}

По формуле ошибки~\eqref{eq:ctd} неприятное событие --- удар, горький вкус --- должно вызывать провал: вышло хуже, чем ожидалось. Часть \nt{DOPA}-нейронов так и делает, а другая отвечает на неприятное \emph{всплеском}, как на награду~\cite{Matsumoto2009}. Эти нейроны сообщают не знак ошибки, а то, что случилось что-то \emph{важное}: неожиданное, сильное, требующее внимания~\cite{BrombergMartin2010}.

Инженеру удобно думать об этом как о двух сигналах. Первый --- знаковая ошибка $\delta$: куда менять веса. Второй --- её модуль $|\delta|$ или новизна события: \emph{насколько сильно} их менять, то есть скорость обучения; после неожиданного события учиться надо быстрее. По смыслу он близок к норадреналиновому коэффициенту усиления из главы~\ref{sec:neurontypes}.

Есть и третий вариант: отдельный провод для отдельной оси. \nt{DOPA}-нейроны, идущие к одной из областей выбора действий, отвечают не на награду, а на новизну и силу раздражителя и нужны, чтобы животное научилось избегать опасного~\cite{Menegas2018}: по проводу идёт не «лучше или хуже», а «опасно или нет».

Что измерено: разные группы \nt{DOPA}-нейронов отвечают на неприятное и на новое по-разному, и разные выходы учат разному. Как именно мозг пользуется сигналом важности --- как скоростью обучения, как сигналом внимания или как отдельной ошибкой по оси опасности, --- обсуждается.

\section{Не только учить, но и подгонять}

У дофамина есть и немедленное действие: чем его больше, тем охотнее и быстрее животное действует. Как совместить это с обучением?

Инженерный ответ --- \emph{разделение по частоте}. Быстрые всплески и провалы длительностью в сотни миллисекунд несут $\delta$ и учат. Медленный уровень, меняющийся за минуты, несёт другую величину --- среднюю награду в единицу времени $\bar R$~\cite{Niv2007}. Пусть действие, выполненное за время $\tau_a$, требует усилия $C/\tau_a$: чем быстрее, тем дороже. Зато каждая секунда промедления стоит $\bar R$ --- столько награды можно было бы за неё получить. Суммарная цена $C/\tau_a+\bar R\,\tau_a$ минимальна при
\begin{empheq}[box=\keybox]{equation}
\tau_a^{*} \;=\; \sqrt{\frac{C}{\bar R}} .
\label{eq:vigor}
\end{empheq}
Чем богаче обстановка, тем дороже время и тем быстрее выгодно действовать: если $\bar R$ удвоилась, время действия уменьшается в $\sqrt2\approx1.4$ раза. Заметьте, что $\bar R$ --- та самая ожидаемая награда, которую вычитали в главе~\ref{sec:dofa}.

Выброс дофамина в месте назначения может меняться и без изменения частоты импульсов: его подстраивают местные сигналы у самих выходов~\cite{Mohebi2019}. Но медленные составляющие есть и в самих импульсах: в опытах следующего раздела плавный рост при приближении к награде виден и в частоте \nt{DOPA}-нейронов~\cite{Kim2020}. Значит, медленный уровень складывается из двух источников --- частоты импульсов и местной регулировки выброса, --- и какой из них главный, по-видимому, зависит от выхода и от задачи.

\begin{keyblock}
\begin{proposition}[Два сигнала на одном проводе]
\label{prop:multiplex}
\nt{DOPA}-система передаёт по меньшей мере две величины, разнесённые по времени: быструю ошибку $\delta$, которая учит, и медленный уровень, который задаёт цену времени~\eqref{eq:vigor} и тем самым скорость действий. Измерено, что быстрая и медленная составляющие ведут себя по-разному; что медленная равна именно $\bar R$, --- гипотеза.
\end{proposition}
\end{keyblock}

\section{Рост при приближении}

Если животное бежит по коридору к далёкой награде, концентрация дофамина в области выбора действий плавно растёт в течение секунд, пока цель приближается~\cite{Hamid2016}. По главе~\ref{sec:dofa} этого быть не должно: всё идёт по плану, и $\delta$ должна быть нулём.

Первое объяснение: этот рост --- не $\delta$, а само ожидание $V$, сигнал «цель близко, стоит постараться» из предыдущего раздела~\cite{Hamid2016}. Второе объяснение сохраняет $\delta$~\cite{Kim2020}. Если ожидание верно, то при горизонте $\tau_\gamma$ на пути к награде $R$ оно растёт как $V=Re^{-(T-t)/\tau_\gamma}$, и по формуле ошибки~\eqref{eq:ctd} $\dd V/\dd t-V/\tau_\gamma=0$. Но пусть издалека ожидание занижено, а по мере приближения оценка уточняется. Тогда ожидание растёт круче, скажем как $V=Re^{-(T-t)/\tau_v}$ с $\tau_v<\tau_\gamma$, и
\begin{empheq}[box=\keybox]{equation}
\delta(t) \;=\; \frac{\dd V}{\dd t}-\frac{V}{\tau_\gamma} \;=\; V\Bigl(\frac{1}{\tau_v}-\frac{1}{\tau_\gamma}\Bigr) \;>\; 0 .
\label{eq:ramp}
\end{empheq}
Ошибка положительна и растёт вместе с $V$ --- тот самый плавный рост (рис.~\ref{fig:ramp}). При $\tau_\gamma=4\,\text{с}$ и $\tau_v=1\,\text{с}$ получается $\delta=0.75\,V$ в секунду. Эту версию проверяли в виртуальном коридоре, мгновенно перенося животное ближе к цели: дофамин отвечал всплеском, как и положено производной, а не плавному уровню~\cite{Kim2020}.

\begin{figure}[t]
\centering
\begin{tikzpicture}[>=Stealth,x=1.6cm,y=2.6cm]
\draw[->,gray!70!black] (-4.2,0) -- (0.5,0) node[below left,black] {\small $t$};
\draw[->,gray!70!black] (-4.2,0) -- (-4.2,1.2);
\foreach \x/\l in {-4/{$T-4$},-2/{$T-2$},0/{$T$}}{\draw[gray!60!black] (\x,-0.03) -- (\x,0.03); \node[below,font=\scriptsize] at (\x,0) {\l\,с};}
\draw[thick,densely dashed,gray!60!black] plot[domain=-4:0,samples=40] (\x,{exp(\x/4)});
\draw[very thick,blue!60!black] plot[domain=-4:0,samples=60] (\x,{exp(\x)});
\draw[very thick,red!70!black] plot[domain=-4:0,samples=60] (\x,{0.75*exp(\x)});
\node[anchor=south east,font=\small,gray!50!black] at (-1.3,0.77) {$V$ при $\tau_v=\tau_\gamma$: $\delta=0$};
\node[right,font=\small,blue!60!black] at (0.05,1.0) {$V$ при $\tau_v=1\,\text{с}$};
\node[right,font=\small,red!70!black] at (0.05,0.72) {$\delta$ по~\eqref{eq:ramp}};
\end{tikzpicture}
\caption{Рост дофамина при приближении к награде как ошибка предсказания, $\tau_\gamma=4\,\text{с}$. Штриховая линия --- ожидание, растущее ровно как требует горизонт ($\delta=0$); синяя --- ожидание, растущее круче; красная --- ошибка~\eqref{eq:ramp}.}
\label{fig:ramp}
\end{figure}

Что измерено: плавный рост есть --- и в концентрации дофамина, и в частоте импульсов \nt{DOPA}-нейронов~\cite{Kim2020}, --- и мгновенные скачки обстановки вызывают скачки дофамина. Какое из двух объяснений верно --- или оба верны на разных выходах, --- обсуждается.

\section{Взгляд назад}

Все теории выше смотрят \emph{вперёд}. Одна из современных теорий переворачивает направление~\cite{Jeong2022}. Пусть мозг учится не предсказывать награду по предвестнику, а, получив награду, \emph{оглядываться}: какое событие чаще других ей предшествовало? Мера такой связи --- разность двух вероятностей:
\begin{empheq}[box=\keybox]{equation}
M \;=\; P(\text{предвестник перед наградой}) \;-\; P(\text{предвестник в случайном окне}) .
\label{eq:retro}
\end{empheq}
Если предвестник часто бывает и без награды, второе слагаемое велико, и связь слаба. В этой теории дофамин сообщает не ошибку предсказания, а то, насколько данное событие --- вероятная \emph{причина} награды.

Механизм оглядки требует того же, что и правило трёх множителей главы~\ref{sec:dofa}: каждое событие должно оставлять тающий след, чтобы в момент награды можно было прочитать, что было перед ней. Разница не в синапсе, а в том, что передаёт \nt{DOPA}.

В опытах, где теория~\eqref{eq:retro} и ошибка~\eqref{eq:ctd} предсказывают разное, выброс дофамина в ряде случаев следовал первой~\cite{Jeong2022}. Но в других опытах ответ дофамина при обучении постепенно сдвигался от момента награды к моменту предвестника --- как требует обучение по ошибке~\eqref{eq:ctd}~\cite{Amo2022}. Какая теория описывает мозг, пока неизвестно.

\section{Ошибка не только в награде}

Последнее расширение --- о том, \emph{о чём} ошибка. Если заменить ожидаемый сок другим, таким же ценным, но иного вкуса, ценность не изменилась, и ошибка~\eqref{eq:ctd} равна нулю. Однако \nt{DOPA}-нейроны на такую замену отвечают~\cite{Takahashi2017}. А искусственно вызванные всплески дофамина могут научить животное связи между двумя нейтральными раздражителями, вообще без награды~\cite{Sharpe2017}. Похоже, \nt{DOPA} сообщает об ошибке предсказания не только награды, но и того, \emph{что} произойдёт.

Формально это простое обобщение~\eqref{eq:ctd}: вместо одного потока награды $r$ берётся набор признаков происходящего $\varphi_k$ --- вкус, место, звук, --- и для каждого своё ожидание $M_k$ и своя ошибка~\cite{Gardner2018}:
\begin{empheq}[box=\keybox]{equation}
\delta_k(t) \;=\; \varphi_k(t) \;+\; \frac{\dd M_k}{\dd t} \;-\; \frac{M_k}{\tau_\gamma} ,
\qquad k=1,\dots,K .
\label{eq:vectortd}
\end{empheq}
Награда --- лишь один из признаков, и вектор ошибок учит не только тому, что хорошо, но и тому, что за чем следует, --- модели мира.

С этим согласуется неоднородность \nt{DOPA}-нейронов: в одной задаче разные нейроны несут разные величины --- одни связаны с наградой, другие с движением, третьи с тем, куда направлено внимание~\cite{Engelhard2019}.

\begin{keyblock}
\begin{proposition}[Жгут вместо провода]
\label{prop:bundle}
Сигнал \nt{DOPA}-системы --- не одно число. Он разложен по нейронам (распределение~\eqref{eq:expectile}, разные признаки~\eqref{eq:vectortd}, отдельная ось опасности) и по времени (быстрая ошибка и медленный уровень~\eqref{eq:vigor}). Скалярная ошибка~\eqref{eq:ctd} --- первое приближение к этому сигналу, как сумматор с порогом --- первое приближение к нейрону.
\end{proposition}
\end{keyblock}

\section{Итог}

\begin{table}[t]
\centering
\footnotesize
\setlength{\tabcolsep}{4pt}
\renewcommand{\arraystretch}{1.15}
\begin{tabular}{>{\raggedright}p{2.5cm}>{\raggedright}p{3.2cm}>{\raggedright}p{3.6cm}>{\raggedright\arraybackslash}p{4.2cm}}
\toprule
Теория & Что идёт по проводу & Главное измерение & Что известно \\
\midrule
ошибка предсказания (гл.~\ref{sec:dofa}) & скаляр $\delta$~\eqref{eq:ctd} & всплеск, перенос на предвестник, провал & измерено; первое приближение \\
распределение & набор $\delta_i$ с разной асимметрией & разные точки переворота у разных нейронов & согласие с опытом есть; роль разброса --- гипотеза \\
важность & $|\delta|$, новизна; ось опасности & всплески на неприятное у части нейронов & измерено; как используется --- обсуждается \\
цена времени & медленный уровень $\bar R$ & дофамин ускоряет действия; медленная составляющая --- и в выбросе у выходов, и в частоте импульсов & разделение быстрого и медленного измерено; $\bar R$ --- гипотеза \\
рост при приближении & $V$ или $\delta$ при уточнении оценки~\eqref{eq:ramp} & плавный рост, скачки при переносе в коридоре & рост измерен; объяснение обсуждается \\
взгляд назад & мера причинной связи~\eqref{eq:retro} & опыты, где предсказания двух теорий расходятся & результаты противоречат друг другу \\
вектор ошибок & $\delta_k$ по признакам~\eqref{eq:vectortd} & ответ на смену вкуса; обучение без награды & измерено; векторная форма --- гипотеза \\
\bottomrule
\end{tabular}
\caption{Что передаёт \nt{DOPA}-система: стандартная картина главы~\ref{sec:dofa} и её современные расширения.}
\label{tab:dofaalt}
\end{table}

Теории таблицы~\ref{tab:dofaalt} не исключают друг друга: почти все сохраняют ошибку предсказания как основу и расширяют её формат. Всерьёз конкурирует с ней только взгляд назад, и этот спор решат опыты. Для инженера вывод простой: цена широковещания из утверждения~\ref{prop:broadcastcost} падает, если провод на самом деле --- много проводов с разными адресатами.

\section{Задачи}

\begin{exercise}[\lvlA]
\label{ex:07:1}
\emph{Точки распределения для одной капли.} Сок размера $1$ приходит с вероятностью $p=0.2$, иначе награда $0$. Найдите по~\eqref{eq:expectile} $V_i$ для $\kappa_i=0.2$, $0.5$, $0.8$. Чему равна медиана этой награды, и чем точка~\eqref{eq:expectile} отличается от квантиля?
\end{exercise}

\begin{exercise}[\lvlB]
\label{ex:07:2}
\emph{Распределение по двум нейронам.} Награда равна $0$ или $x$, вероятность $x$ равна $p$. Два нейрона с правилом~\eqref{eq:asymmetric} сообщили $V(1/2)=0.25$ и $V(0.8)=0.5$. Восстановите $p$, $x$ и стандартное отклонение награды. Что сообщит нейрон с $\kappa=0.2$?
\end{exercise}

\begin{exercise}[\lvlB]
\label{ex:07:3}
\emph{Моделирование: распределение из асимметрии.} Смоделируйте девять \nt{DOPA}-нейронов с $\kappa_i=0.1,\dots,0.9$ и правилом~\eqref{eq:asymmetric}, $\alpha_i^+=2\kappa_i$, $\alpha_i^-=2(1-\kappa_i)$, при награде, равномерной на $[0,1]$. Как разброс $V_i$ зависит от $\eta$, и какое $\eta$ нужно, чтобы отличить это распределение от двух равновероятных капель $0.25$ и $0.75$?
\end{exercise}

\begin{exercise}[\lvlA]
\label{ex:07:4}
\emph{Цена времени.} (а)~Выведите~\eqref{eq:vigor} и найдите суммарную цену в оптимуме. (б)~Мозг оценил $\bar R$ с ошибкой в $k$ раз. Найдите относительную переплату при $k=2$ и $k=4$. Насколько точно медленный уровень дофамина должен сообщать $\bar R$?
\end{exercise}

\begin{exercise}[\lvlB]
\label{ex:07:5}
\emph{Всплеск при переносе.} Ожидание растёт по~\eqref{eq:ramp} с $\tau_v=1\,\text{с}$, $\tau_\gamma=4\,\text{с}$; животное мгновенно переносят из точки $T-2\,\text{с}$ в $T-1\,\text{с}$. Найдите площадь всплеска $\delta$ в долях $R$ и сколько секунд рампы до переноса он заменяет. Что покажет перенос, если дофамин сообщает само $V$?
\end{exercise}

\begin{exercise}[\lvlB]
\label{ex:07:6}
\emph{Проектирование: два сигнала на проводе.} По проводу \nt{DOPA} идут уровень, растущий до $s=1/60$ импульса в секунду за секунду, и события по $A=3$ импульса. Синапс со следом $\tau_e=1\,\text{с}$ вычитает из частоты её копию, сглаженную с $\tau_f$. При каких $\tau_f$ теряется не больше $10\,\%$ события, а вклад уровня за время следа меньше $10\,\%$ вклада события?
\end{exercise}

\begin{exercise}[\lvlB]
\label{ex:07:7}
\emph{Проектирование опыта: взгляд назад против ошибки.} В окне предвестник бывает с вероятностью $0.1$, после него награда --- с вероятностью $0.8$. Сравните $\Delta P=P(\text{нагр.}\mid\text{предв.})-P(\text{нагр.}\mid\text{нет})$ и $M$ из~\eqref{eq:retro}, если (а)~добавить награды без предвестника, $0.08$ на окно; (б)~удвоить частоту предвестника без наград.
\end{exercise}

\begin{exercise}[\lvlC]
\label{ex:07:8}
\emph{Жгут и цена широковещания.} В модели задачи~\ref{ex:06:8} $N$ нейронов разбиты на $K$ групп со своими проводами, и в каждый провод просачивается доля $\rho$ чужих ошибок. Покажите, что постоянная обучения равна $N/K+2+\rho^2N(1-1/K)$. Сколько попыток она даёт при $K\to\infty$, $N=10^{10}$ и $\rho=0.01$?
\end{exercise}
